\section{Sturmian discrepancy estimates}\label{sec:fp:tv}

The rotation model gives deterministic estimates for the finite-window laws considered above. The argument combines standard discrepancy bounds with deterministic pushforward under $\Phi_{m,n}$.

\subsection{Stabilized block laws from discrepancy}

Fix a stabilization aperture $m$ and a stabilized block length $n\ge 1$. Here $\Phi_{m,n}$ denotes the finite-block sliding map
\[
\Phi_{m,n}(a_0\ldots a_{m+n-2})
:=
\bigl(\Fold_m(a_t\ldots a_{t+m-1})\bigr)_{t=0}^{n-1}.
\]
For the first $N$ starting positions $\{x_0+t\alpha\}_{t=0}^{N-1}$ under the standard two-interval rotation coding, let $\hat{\mu}_{m+n-1,N}$ denote the empirical distribution of raw words of length $m+n-1$ beginning at those positions, and let $\mu_{m+n-1}$ be the corresponding limiting raw-word distribution obtained from Lebesgue measure on the depth-$(m+n-1)$ Sturmian atoms. Define the empirical stabilized $n$-block law and its limit by
\[
\hat{\eta}_{m,n,N}:=(\Phi_{m,n})_*\hat{\mu}_{m+n-1,N},
\qquad
\eta_{m,n}:=(\Phi_{m,n})_*\mu_{m+n-1}.
\]
For $n=1$ these are exactly the empirical folded histogram and its limit:
\[
\hat{\eta}_{m,1,N}=\hat{\pi}_{m,N},
\qquad
\eta_{m,1}=\pi_m^{\mathrm{fold}}.
\]
Let
\[
\mathcal Y_{m,n}:=\Phi_{m,n}(\{0,1\}^{m+n-1})
\]
denote the finite support of these stabilized $n$-block laws.

\begin{theorem}[Stabilized block law from discrepancy]\label{thm:tv-certificate}
Let $D_N^*(\alpha;x_0)$ denote the interval discrepancy of the orbit segment
$\{x_0+t\alpha\}_{t=0}^{N-1}$, namely the supremum over circle intervals $I$
of
\[
\left|
\frac1N\#\{0\le t<N:\{x_0+t\alpha\}\in I\}-|I|
\right|.
\]
Then, for every $m\ge2$ and $n\ge1$,
\[
D_{\TV}\bigl(\hat{\eta}_{m,n,N},\eta_{m,n}\bigr)
\le (m+n)D_N^*(\alpha;x_0).
\]
This is a supporting rotation-source certificate; the structural conjugacy
results above do not depend on it.
\end{theorem}

\begin{proof}
Put $L=m+n-1$.  In the standard two-interval Sturmian coding, the depth-$L$
raw-word atoms form a partition of the circle into at most $L+1=m+n$
intervals.  For each such interval $I$ the definition of interval discrepancy
gives
\[
\left|
\frac1N\#\{0\le t<N:\{x_0+t\alpha\}\in I\}-|I|
\right|
\le D_N^*(\alpha;x_0),
\]
where $|I|$ denotes Lebesgue length.  Summing over the at most $m+n$ atoms gives
\[
D_{\TV}\bigl(\hat{\mu}_{L,N},\mu_L\bigr)
\le (m+n)D_N^*(\alpha;x_0).
\]
Total variation cannot increase under deterministic pushforward.  Applying
the map $\Phi_{m,n}$ to the two raw-word laws gives the displayed bound.
\end{proof}

\begin{corollary}[Bound for bounded stabilized-block observables]\label{cor:observable-certificate}
If $g:\mathcal Y_{m,n}\to\RR$ is bounded, then
\[
\left|
\int g\,d\hat{\eta}_{m,n,N}-\int g\,d\eta_{m,n}
\right|
\le
2\|g\|_\infty (m+n)D_N^*(\alpha;x_0).
\]
\end{corollary}

\begin{proof}
For probability measures $\lambda,\lambda'$ on a finite set,
$|\int g\,d\lambda-\int g\,d\lambda'|
\le 2\|g\|_\infty D_{\TV}(\lambda,\lambda')$.  Apply
Theorem~\ref{thm:tv-certificate}.
\end{proof}

\subsection{Boundary-layer stability for stabilized blocks}\label{sec:fp:tv:boundary}

Let $R_0:H\to\{0,1\}$ denote the sharp indicator readout and let $R_\rho:H\to\{0,1\}$ be a binary finite-resolution readout obtained from a smoothed kernel with resolution parameter $\rho$. Write
\[
\varepsilon_\rho
:=
\mu\bigl(\{h\in H:\ R_\rho(h)\neq R_0(h)\}\bigr)
\]
for the boundary-layer mass. Let $\mu_{m+n-1}^{(0)}$ and $\mu_{m+n-1}^{(\rho)}$ be the corresponding raw-word marginals on $\{0,1\}^{m+n-1}$, and let
\[
\eta_{m,n}^{(0)}:=(\Phi_{m,n})_*\mu_{m+n-1}^{(0)},
\qquad
\eta_{m,n}^{(\rho)}:=(\Phi_{m,n})_*\mu_{m+n-1}^{(\rho)}
\]
be the stabilized $n$-block laws.

\begin{proposition}[Boundary-layer control for raw and stabilized block laws]\label{prop:boundary-layer}
With $L=m+n-1$,
\[
D_{\TV}\bigl(\mu_L^{(\rho)},\mu_L^{(0)}\bigr)
\le L\varepsilon_\rho
\]
and therefore
\[
D_{\TV}\bigl(\eta_{m,n}^{(\rho)},\eta_{m,n}^{(0)}\bigr)
\le (m+n-1)\varepsilon_\rho .
\]
\end{proposition}

\begin{proof}
Couple the sharp and blurred readouts by evaluating both on the same source
point and along the same orbit segment.  The two length-$L$ raw words can
differ only if at least one of the $L$ sampled readout symbols lies in the
boundary-error set $\{h:R_\rho(h)\ne R_0(h)\}$.  By the union bound, this event
has probability at most $L\varepsilon_\rho$.  The coupling characterization of
total variation gives the raw-word bound.  The stabilized law is the
deterministic pushforward of the raw-word law under $\Phi_{m,n}$, so data
processing gives the second bound.
\end{proof}

\begin{corollary}[Interval-partition resolution certificate]\label{cor:resolution-certificate}
The blurred readout need not be a two-interval coding. Assume instead that its
depth-$(m+n-1)$ atom partition consists of at most $m+n$ intervals, so that its
empirical raw-word law differs from its limiting raw-word law in total
variation by at most $(m+n)D_N^*(\alpha;x_0)$. Then, for the empirical blurred
stabilized law $\hat{\eta}_{m,n,N}^{(\rho)}$,
\[
D_{\TV}\bigl(\hat{\eta}_{m,n,N}^{(\rho)},\eta_{m,n}^{(0)}\bigr)
\le
(m+n)D_N^*(\alpha;x_0)+(m+n-1)\varepsilon_\rho.
\]
\end{corollary}

\begin{remark}
For the sharp two-interval Sturmian coding, the depth-$(m+n-1)$ atoms form at
most $m+n$ intervals, so the same argument gives the displayed certificate
with $\varepsilon_\rho=0$.  The present statement records the
interval-partition extension: the certificate holds for any blurred readout
whose depth atoms form at most $m+n$ intervals, even when the blurred readout
is not itself a two-interval coding.
\end{remark}

\begin{proof}
The stated interval-complexity hypothesis gives the same empirical-to-limiting
bound for the blurred raw-word law as in Theorem~\ref{thm:tv-certificate}, and
data processing gives the corresponding stabilized-law bound. Combine this with
Proposition~\ref{prop:boundary-layer} and apply the triangle inequality.
\end{proof}

\subsection{Recovery of finite-range observables}

\begin{theorem}[Recovery of finite-range observables from stabilized blocks]\label{thm:operational-observable-recovery}
Assume $m\ge3$ and $n\ge m$.  Let
$f:\{0,1\}^{m+n-1}\to\RR$ be any finite-range raw observable.  There is a
unique observable $\widetilde f:\mathcal Y_{m,n}\to\RR$ such that
\[
f=\widetilde f\circ\Phi_{m,n}
\quad\text{on }\{0,1\}^{m+n-1}.
\]
It is given by $\widetilde f(y)=f(\Psi_{m,n}(y))$, where
$\Psi_{m,n}$ is the inverse block decoder from
Theorem~\ref{thm:block-invertibility}.  Consequently
\[
\int f\,d\hat{\mu}_{m+n-1,N}
=
\int \widetilde f\,d\hat{\eta}_{m,n,N},
\qquad
\int f\,d\mu_{m+n-1}
=
\int \widetilde f\,d\eta_{m,n},
\]
and
\[
\left|
\int f\,d\hat{\mu}_{m+n-1,N}
-
\int f\,d\mu_{m+n-1}
\right|
\le
2\|f\|_\infty (m+n)D_N^*(\alpha;x_0).
\]
\end{theorem}

\begin{proof}
For $m\ge3$ and $n\ge m$, Theorem~\ref{thm:block-invertibility} makes
$\Phi_{m,n}$ a bijection from $\{0,1\}^{m+n-1}$ onto
$\mathcal Y_{m,n}$, with inverse $\Psi_{m,n}$.  Therefore the formula
$\widetilde f(y)=f(\Psi_{m,n}(y))$ is the unique function satisfying
$f=\widetilde f\circ\Phi_{m,n}$.  The two identities of integrals are simply
change of variables under deterministic pushforward.  Since
$\|\widetilde f\|_\infty=\|f\|_\infty$, the last estimate is
Corollary~\ref{cor:observable-certificate} applied to $\widetilde f$.
\end{proof}

\begin{corollary}[Bounded-type slopes]\label{cor:tv-bounded-type}
If $\alpha$ has bounded continued-fraction partial quotients, then there is a
constant $C_\alpha$ such that, for all $N\ge2$,
\[
D_{\TV}\bigl(\hat{\eta}_{m,n,N},\eta_{m,n}\bigr)
\le
C_\alpha(m+n)\frac{\log N}{N}.
\]
For $m\ge3$ and $n\ge m$, every recovered finite-range observable in
Theorem~\ref{thm:operational-observable-recovery} satisfies the corresponding
bound
\[
\left|
\int f\,d\hat{\mu}_{m+n-1,N}
-
\int f\,d\mu_{m+n-1}
\right|
\le
2C_\alpha\|f\|_\infty(m+n)\frac{\log N}{N}.
\]
\end{corollary}

\begin{proof}
The classical discrepancy estimate for irrational rotations with bounded
partial quotients gives
$D_N^*(\alpha;x_0)\le C_\alpha\log N/N$, uniformly in the starting point
\cite{KuipersNiederreiter1974,DrmotaTichy1997}.  Substitute this estimate into
Theorem~\ref{thm:tv-certificate} and
Theorem~\ref{thm:operational-observable-recovery}.
\end{proof}

\begin{remark}
Theorem~\ref{thm:operational-observable-recovery} and Corollary~\ref{cor:resolution-certificate} quantify the rotation source law after the conjugacy has been established. The finite-sample error terms are the orbit discrepancy and the boundary-layer mass.
\end{remark}
